\chapter{The local threefold benchmark}
\label{cylocal:ch:benchmark}

The affine threefold $\mathbb C^3$ gives a concrete setting in which to
compare Hall operations with fields on a curve. Its category of
finite-length sheaves has a critical Hall algebra. A free holomorphic
gauge theory on the same threefold has curve defects and their
operator product expansions (OPEs). We construct these two inputs
with their own coefficients.

The Hall algebra admits an affine-Yangian
presentation over its polynomial coefficient ring. The
defect has an explicit Gaussian Ward identity. A comparison between
them must preserve the operations computed on each side.

\section{A potential and its critical coefficients}
\label{cylocal:sec:critical}

Let $Q$ have one vertex and three loops $x,y,z$. Its path algebra is
$\mathbb C\langle x,y,z\rangle$. Identify cyclic rotations in the
vector space of words, and take the potential
\begin{equation}\label{cylocal:eq:potential}
 W=zxy-zyx=z[x,y].
\end{equation}

The cyclic derivative with respect to a generator sums over its
occurrences, removing each occurrence and reading the remaining
word from the next position. Thus
\[
 \partial_xW=yz-zy,\qquad
 \partial_yW=zx-xz,\qquad
 \partial_zW=xy-yx.
\]
The Jacobi algebra, obtained by imposing these three derivatives,
is $\mathbb C[x,y,z]$. Its finite-dimensional modules are precisely
finite-length coherent sheaves on $\mathbb C^3$.

Let $T=(\mathbb C^*)^2$ act on the loops with additive characters
$a,b,-a-b$. Each term of $W$ has weight zero. Put
\[
 R=H_T^\bullet(\mathrm{pt},\mathbb Q)=\mathbb Q[a,b],
 \qquad |a|=|b|=2.
\]

Charge $n\geq0$ means representation dimension $n$. The smooth
representation stack and its function are
\[
 \mathcal M_n=[\operatorname{Mat}_n(\mathbb C)^3/GL_n],
 \qquad f_n(X,Y,Z)=\operatorname{Tr}Z[X,Y].
\]
The classical critical locus parametrizes commuting triples. The coefficient
complex also records the singularity of $f_n$ in its smooth ambient
stack. It is additional to the underlying set of commuting triples.

Write $\Phi_f$ for shifted vanishing cycles, extended from the zero
fiber, with convention $\Phi_0\mathbb Q=\mathbb Q$. The virtual
shift on $\mathcal M_n$ is its complex dimension $2n^2$. We use the
relative equivariant convention: quotienting by the external torus
does not subtract $\dim T$ from that shift. Define
\begin{equation}\label{cylocal:eq:critical-hall}
 \mathcal H_n^T=
 H^\bullet(\mathcal M_n/T,\Phi_{f_n}\mathbb Q[2n^2]),
 \qquad \mathcal H^T=\bigoplus_{n\geq0}\mathcal H_n^T.
\end{equation}

The extension correspondence, with its critical coefficient maps,
gives the Hall product. It adds charge and preserves cohomological
degree. Its unit is $1\in\mathcal H_0^T=R$.

At charge one the function is zero, so
\begin{equation}\label{cylocal:eq:charge-one}
 \mathcal H_1^T=R[u][2],\qquad |u^r|=2r-2.
\end{equation}
Here $u=c_1(L)$ for the universal one-dimensional representation $L$.
It is independent of the external parameters $a,b$.
We fix the root convention for the Hall product below.
In particular, charge one has infinite rank over $R$.

\section{The cut and the product sign}
\label{cylocal:sec:cut}

The potential is linear in every entry of $Z$. Forgetting $Z$
therefore leads to the moment equation $[X,Y]=0$. Let
\[
 C_n=\{(X,Y)\in\operatorname{Mat}_n(\mathbb C)^2:[X,Y]=0\},
 \qquad \mathcal P_n^T=H^{\mathrm{BM},T}_\bullet([C_n/GL_n],\mathbb Q).
\]
We regard Borel--Moore homology as cohomologically graded by reversing
its index. Its equivariant normalization has
$\mathcal P_0^T=R$ and $\mathcal P_1^T=R[u][2]$.
The product on $\mathcal P^T=\bigoplus_n\mathcal P_n^T$ is the
refined pullback along the moment correspondence followed by the
proper map forgetting the invariant subspace.

The refinement already matters for a two-dimensional representation.
For an invariant first coordinate, write
\[
 X=\begin{pmatrix}x_1&r\\0&x_2\end{pmatrix},\qquad
 Y=\begin{pmatrix}y_1&v\\0&y_2\end{pmatrix}.
\]
Then
\begin{equation}\label{cylocal:eq:residual}
 [X,Y]=\begin{pmatrix}0&\rho\\0&0\end{pmatrix},
 \qquad \rho=(x_1-x_2)v-(y_1-y_2)r.
\end{equation}

Each diagonal block commutes, but the extension must still satisfy
$\rho=0$. On the coincident diagonal $x_1=x_2$, $y_1=y_2$, this
equation vanishes identically. Its derived zero fiber retains the
odd generator of the Koszul complex
\[
 \mathbb Q[x_1,x_2,y_1,y_2,r,v][\eta],
 \qquad |\eta|=-1,\quad d\eta=\rho.
\]
After derived specialization to that diagonal, $d\eta=0$.
Ordinary restriction of the classical zero locus would lose this
negative-degree class.

Dimensional reduction provides isomorphisms of graded $R$-modules
$\theta_n:\mathcal P_n^T\to\mathcal H_n^T$. They arise by
integrating every cut-loop entry, so the equations imposed on a flag
include \eqref{cylocal:eq:residual}. With complex orientations, the
two ordinary Hall products obey
\begin{equation}\label{cylocal:eq:product-sign}
 m_{\mathrm{crit}}(\theta_n\otimes\theta_m)
   =(-1)^{nm}\theta_{n+m}m_{\mathrm{red}}.
\end{equation}
This is the product comparison in
\cite{cylocal:RS}*{Appendix, Theorem 4.4} and
\cite{cylocal:Davison}*{\S3.9}.

The sign has a two-variable origin. For the quadratic potential
$st$ on $\mathbb A^2$, the compact critical class maps to
$(1,-1)$ in the two complex-oriented axes of $\{st=0\}$.
Both axes specialize with multiplicity one to the smooth fiber,
so their sum map has kernel generated by $(1,-1)$.
Restriction and integration along the two axes consequently differ
by $-1$.

The flag comparison has $nm$ such paired normal directions.
Thom--Sebastiani multiplies their signs and gives
\eqref{cylocal:eq:product-sign}. The remaining maps are the same
proper flag map and the same refined moment pullback. Thus this
comparison is defined before Euler classes are inverted.

\begin{proposition}\label{cylocal:prop:integral-reduction}
The maps
\[
 \widehat\theta_n=(-1)^{\binom n2}\theta_n
\]
give a unital isomorphism of charge-graded, cohomologically graded
$R$-algebras $\mathcal P^T\simeq\mathcal H^T$.
\end{proposition}
\begin{proof}
The identity
$\binom{n+m}{2}=\binom n2+\binom m2+nm$ absorbs exactly the sign
in \eqref{cylocal:eq:product-sign}. At charge zero both reduction
and its scalar correction are the identity. Each $\theta_n$ is an
isomorphism by linear dimensional reduction, so their direct sum
has an inverse. All scalars used are rational units.
\end{proof}

This theorem concerns the rational Hall operations. A comparison
of full derived flag coefficient categories must also retain the
closed odd directions illustrated by \eqref{cylocal:eq:residual}.

\section{The full algebra and its charge-one generators}
\label{cylocal:sec:yangian}

The spherical subalgebra of $\mathcal H^T$ is the unital
$R$-subalgebra generated by $\mathcal H_1^T$. In the present model
it is the whole algebra. We impose no nilpotence condition on any
of the three matrices.

Put $a_1=a$, $a_2=b$, $a_3=-a-b$ and
\[
 \sigma_2=a_1a_2+a_1a_3+a_2a_3,
 \qquad \sigma_3=a_1a_2a_3=-ab(a+b).
\]
Define $Y_R^+$ to be the unital associative $R$-algebra with
generators $e_r$, $r\geq0$, of charge one and degree $2r-2$.
Its relations are
\begin{align}
 &[e_{r+3},e_s]-3[e_{r+2},e_{s+1}]
     +3[e_{r+1},e_{s+2}]-[e_r,e_{s+3}]\notag\\
 &\hspace{2em}
   +\sigma_2([e_{r+1},e_s]-[e_r,e_{s+1}])
       =-\sigma_3(e_re_s+e_se_r),
       \label{cylocal:eq:yangian-quadratic}\\
 &\sum_{\pi\in S_3}
   [e_{r_{\pi(1)}},[e_{r_{\pi(2)}},e_{r_{\pi(3)}+1}]]=0.
       \label{cylocal:eq:yangian-serre}
\end{align}
All indices are nonnegative. This is the strictly positive affine
Yangian presentation. No Cartan or negative generators occur in
its definition.

\begin{theorem}[Davison]\label{cylocal:thm:yangian}
The assignment of $e_r$ to the normalized charge-one class $u^r$
extends to an isomorphism
\[
 Y_R^+\xrightarrow{\ \cong\ }\mathcal H^T.
\]
It preserves charge, cohomological degree and the unit. It requires
neither Euler-class localization nor completion in charge.
Let any torus act diagonally on the three loops and preserve $W$.
Its full critical Hall algebra is generated over its equivariant
coefficient ring by charge one.
\end{theorem}
\begin{proof}
The algebra comparison is \cite{cylocal:Davison}*{Theorem 5.11},
using the rational dimensional-reduction normalization of \S3.9
there. Its generators are the powers of the determinant class in
charge one. The defining relations are (40)--(41) of that source,
with the root convention made explicit in
\eqref{cylocal:eq:shuffle} below. The all-torus generation statement
is its Theorem 5.6. The hypotheses hold here because each cyclic
term of \eqref{cylocal:eq:potential} has total torus weight zero.
\end{proof}

The coefficient convention is essential for supported variants.
At charge one, the cut-loop nilpotent locus is $z=0$ in
$\mathbb A^3\times B\mathbb G_m$. Ordinary restriction of the critical
coefficient $\mathbb Q[2]$ gives $\mathbb Q[2]$ there. Extraordinary
restriction gives $\mathbb Q[0](-1)$, since the normal complex line
contributes $[-2](-1)$. These define different coefficient problems.
The Tate twist $(-1)$ records the normal line's Hodge weight when
that refinement is retained.

Theorem~\ref{cylocal:thm:yangian} concerns the unrestricted stack
and its displayed coefficient. A supported or nilpotent sector
requires its own restriction functor, shift and generation theorem.

\section{A class lost by specialization of the shuffle map}
\label{cylocal:sec:specialization}

The full-rank torus gives a faithful polynomial realization
\cite{cylocal:Davison}*{Proposition 3.19}. Its ambient charge-$n$
module is
\[
 \mathcal A_n=R[x_1,\ldots,x_n]^{S_n}[2n^2].
\]
Thus a polynomial of ordinary cohomological degree $d$ has adjusted
degree $d-2n^2$. Put $\mathcal A=\bigoplus_{n\geq0}\mathcal A_n$
and define
\[
 \zeta(v)=\frac{(v+a)(v+b)(v-a-b)}{v}
        =v^2+\sigma_2+\frac{\sigma_3}{v}.
\]

We orient a cross-block difference as second root minus first root.
For $f\in\mathcal A_n$ and $g\in\mathcal A_m$, the product is
\begin{equation}\label{cylocal:eq:shuffle}
 (f\star g)(x)=
 \sum_{\substack{I\sqcup J=\{1,\ldots,n+m\}\\|I|=n}}
 f(x_I)g(x_J)\prod_{i\in I,\,j\in J}\zeta(x_j-x_i).
\end{equation}
Residues cancel in pairs across $x_i=x_j$, so this expression is
polynomial. The three-block product contains one factor for each
ordered cross-block pair, which proves associativity.

The convention with $\zeta(x_i-x_j)$ is carried to
\eqref{cylocal:eq:shuffle} by the $R$-linear map
$f(x_1,\ldots,x_n)\mapsto f(-x_1,\ldots,-x_n)$.
The same map sends the charge-one monomial $x^r$ to $(-x)^r$.
Reversing the root direction therefore requires this simultaneous
change of generator coordinates. With \eqref{cylocal:eq:shuffle},
the normalized Hall map is
\begin{equation}\label{cylocal:eq:ambient-map}
 \Phi_T:\mathcal H^T\hookrightarrow\mathcal A,
 \qquad \Phi_T(e_r)=x^r.
\end{equation}
Here $e_r$ also denotes its image in $\mathcal H^T$ under
Theorem~\ref{cylocal:thm:yangian}.

For two roots $x,y$, put $v=y-x$. Direct multiplication gives
\begin{align*}
 \Phi_T([e_1,e_0])
  &=(x-y)\bigl(\zeta(v)-\zeta(-v)\bigr)\\
  &=-2\sigma_3=2ab(a+b).
\end{align*}
The result is a constant polynomial in charge two. Its adjusted
degree is $6-8=-2$, which agrees with the degree of the commutator.
It is not an element of the charge-zero coefficient ring.

Set $b=0$. The specialized kernel is $v^2-a^2$, an even function.
Exchanging two blocks in \eqref{cylocal:eq:shuffle} now leaves every
factor unchanged, so the entire specialized ambient algebra is
commutative. The Hall algebra at the same parameter is still
noncommutative. An explicit model exhibits the surviving class.

Let $s$ be a coordinate on $\mathbb C^*$ and let
$D=s\,\mathrm d/\mathrm ds$. Differential operators satisfy
\begin{equation}\label{cylocal:eq:operator-product}
 (s^nP(D))(s^mQ(D))=s^{n+m}P(D+m)Q(D).
\end{equation}
Inside the Lie algebra of such operators with Laurent parameter
$a$, take the free $\mathbb Q[a]$-module
\begin{equation}\label{cylocal:eq:rees}
 \mathfrak w_a=
 \bigoplus_{n\geq1,\ r\geq0}
       \mathbb Q[a]\,s^nD^r a^{r-1}.
\end{equation}
Its displayed generators form a basis, with charge $n$ and degree
$2r-2$. The symbol $s^n/a$ is a basis
element of this module. Scalar multiplication by $a^{-1}$ is not
allowed.

Formula~\eqref{cylocal:eq:operator-product} proves closure under
commutators. The highest powers of $D$ cancel. Each remaining term
has a nonnegative coefficient power of $a$ in the displayed basis.

The one-parameter comparison is
\begin{equation}\label{cylocal:eq:one-parameter}
 \mathcal H^T\otimes_R\mathbb Q[a]
   \cong U_{\mathbb Q[a]}(\mathfrak w_a),
 \qquad e_r\longmapsto sD^r a^{r-1},
\end{equation}
where $R\to\mathbb Q[a]$ sends $b$ to zero
\cite{cylocal:Davison}*{Theorem 5.9}.
The Hall base change itself is justified by equivariant purity
\cite{cylocal:Davison}*{Proposition 3.6}.
In this model
\begin{equation}\label{cylocal:eq:surviving-commutator}
 [e_1,e_0]\longmapsto[sD,s/a]=s^2/a.
\end{equation}
The right side is a nonzero basis vector of $\mathfrak w_a$, and
remains nonzero in its enveloping algebra by the
Poincar\'e--Birkhoff--Witt theorem.

\begin{proposition}\label{cylocal:prop:specialization-kernel}
The specialization of \eqref{cylocal:eq:ambient-map} at $b=0$
is not injective. Its kernel contains the nonzero Hall class
$[e_1,e_0]$. The specialized Hall algebra is nevertheless generated
by charge one.
\end{proposition}
\begin{proof}
Its ambient image $2ab(a+b)$ vanishes at $b=0$, whereas
\eqref{cylocal:eq:surviving-commutator} proves that the source class
survives. The generation statement is the one-parameter case of
Theorem~\ref{cylocal:thm:yangian}.
\end{proof}

At $a=0$ as well, the basis in \eqref{cylocal:eq:rees} gives symbols
$w_{n,r}$ with bracket
\begin{equation}\label{cylocal:eq:symbol-bracket}
 [w_{n,r},w_{m,j}]=(rm-jn)w_{n+m,r+j-1}.
\end{equation}
When $r=j=0$, the bracket is zero. This is the Poisson bracket of
the symbols $s^n\xi^r$ with $\{\xi,s\}=s$.
The nonequivariant Hall algebra is its enveloping algebra
\cite{cylocal:Davison}*{Theorem 5.4 and Corollary 5.5}.
In particular, $[w_{1,1},w_{1,0}]=w_{2,0}\ne0$.
Setting the equivariant parameters to zero in the ambient shuffle
formula would discard this bracket.

This symbol Lie algebra is generated by the elements $w_{1,r}$.
Indeed, induction on $n\geq2$ uses
\[
 [w_{1,r+1},w_{n-1,0}]=(r+1)(n-1)w_{n,r}.
\]
The coefficient is a nonzero rational unit. Thus charge-one
generation follows from actual brackets.

To lift it equivariantly,
purity supplies free coefficient modules with their nonequivariant
specialization. In each charge their cohomological grading is
bounded below. Lifting products of charge-one generators and
inducting on degree then proves generation over the parameter
ring. This is the graded argument used in
\cite{cylocal:Davison}*{Theorem 5.6}.

\section{An independent free bulk and defect}
\label{cylocal:sec:free-bulk}

We now specify the field-theoretic input. On
$M=\mathbb C_z\times\mathbb C^2_{w_1,w_2}$, with volume form
$\Omega=\mathrm dz\,\mathrm dw_1\,\mathrm dw_2$, take a
one-dimensional abelian gauge algebra with pairing $1$.
The Batalin--Vilkovisky (BV) field complex of free holomorphic
Chern--Simons (hCS) theory is
\[
 \mathcal E=\Omega^{0,\bullet}(M)[1],\qquad Q=\bar\partial.
\]

A tangent component of antiholomorphic form degree $r$ has degree
$r-1$. Its coordinate has ghost degree $1-r$.
Integration of compactly supported variations against $\Omega$
therefore gives a pairing of degree $-1$. The quadratic action is
\begin{equation}\label{cylocal:eq:bulk-action}
 S_0(\mathsf A)=\frac12\int_M\Omega\,
                   \mathsf A\bar\partial\mathsf A.
\end{equation}
The action is a local functional, also defined on compactly
supported fields. There is no cubic interaction for this abelian algebra.

Choose the flat Dolbeault gauge and its Gaussian heat regularization.
The physical component $A=A_{\bar z}$ has ghost number zero.
Its gauge differential, denoted by $s$, acts on $A$ and its ghost $c$ by
\[
 sA=\bar\partial c,\qquad sc=0.
\]

At a formal point, the Dolbeault jet contraction leaves holomorphic
ghost jets. The antiholomorphic Euler contraction divided by total
antiholomorphic polynomial and form degree contracts every positive
degree. This defines the contraction coefficient by coefficient
on formal series.

For nonnegative integers $\alpha_1,\alpha_2,r$,
the resulting ghost algebra has generators
\[
 c_{\alpha,r}=\frac{1}{\alpha_1!\alpha_2!}
  \partial_z^r\partial_{w_1}^{\alpha_1}
                    \partial_{w_2}^{\alpha_2}c,
 \qquad |c_{\alpha,r}|=1.
\]
They anticommute. Their differential is zero and their OPE is regular.
Their normal weight is $|\alpha|=\alpha_1+\alpha_2$.
This is the formal ghost-jet object of the specified free theory.

Place an even $\beta\gamma$ system on the curve
$w_1=w_2=0$. Introduce independent parameters
$\hbar,q$ and $\kappa_{\mathrm{def}}$, with coefficient ring
$\mathbb C[\kappa_{\mathrm{def}}][[\hbar,q]]$.
All three parameters have ghost degree zero.
The covariance is fixed by
\begin{equation}\label{cylocal:eq:covariance}
 \beta(z)\gamma(w)\sim\frac{\hbar}{z-w},\qquad
 \gamma(z)\beta(w)\sim-\frac{\hbar}{z-w},\qquad
 \beta\beta\sim\gamma\gamma\sim0.
\end{equation}

Equivalently, use the Cauchy inverse for the kinetic operator
$-\bar\partial$ with the integration convention
\[
 \tau(f)=\frac{1}{2\pi i}\int_{\mathbb C}
                 \mathrm dz\wedge\mathrm d\bar z\,f.
\]
The state module is the polynomial vacuum module generated by the
fields and their derivatives, completed in the stated parameters.
Its vacuum is denoted by $\mathbf1$.

Write $\partial=\partial_z$ and use primes for this derivative.
Choose the current and its level to be
\begin{equation}\label{cylocal:eq:defect-current}
 J=\beta+\frac{\kappa_{\mathrm{def}}}{2}\partial\gamma,
 \qquad k_{\mathrm{def}}=\hbar\kappa_{\mathrm{def}}.
\end{equation}
Differentiating the two different slots in
\eqref{cylocal:eq:covariance} gives two positive cross contractions.
Each contributes $\hbar\kappa_{\mathrm{def}}/2$, so
\begin{equation}\label{cylocal:eq:defect-level}
 J(z)J(w)\sim\frac{k_{\mathrm{def}}}{(z-w)^2}.
\end{equation}
The level is determined by this covariance and current.

Couple the zeroth normal jet of $A$ by the source exponential
$\exp(q\tau(AJ))$. Sources and ghosts are smooth and compactly
supported in $z$. Higher normal jets may be included formally, with
finite tangential derivative order at each fixed normal weight and
parameter precision. These are the free hCS current couplings of
\cite{cylocal:CP}*{\S6.3.1}, with a specified Gaussian defect and
normalization. The parameter $q$ is the coefficient in this source
exponential. No inverse of $\hbar$ has been taken.

Two such source vertices cannot be joined by a bulk propagator.
Write $\epsilon_{ij\ell}$ for the alternating three-tensor and
$p_\ell$ for the holomorphic momentum components.
The flat physical one-form covariance has a scalar heat multiplier
times $\epsilon_{ij\ell}p_\ell$. It is antisymmetric in its
component indices. Its $\bar z,\bar z$ component is
therefore zero at every cutoff, also after normal differentiation.
The remaining connected source calculation uses only the defect
covariance \eqref{cylocal:eq:defect-level}.

\section{The Ward class and a nonzero regular coupling}
\label{cylocal:sec:ward}

Let $C(u)=1/u$ be the locally integrable Cauchy distribution and
define $u^{-2}=-\partial_uC$. Normalize $\delta^\tau$ by
$\tau(\delta^\tau f)=f(0)$. Since
$\mathrm dz\wedge\mathrm d\bar z=-2i\,\mathrm dx\wedge\mathrm dy$,
the distribution identities are
\begin{equation}\label{cylocal:eq:cauchy}
 \bar\partial C=-\delta^\tau,
 \qquad \bar\partial(u^{-2})=\partial_u\delta^\tau.
\end{equation}
These identities fix the sign in the Ward calculation.

\begin{proposition}\label{cylocal:prop:ward}
For the normalized connected Gaussian source functional
\[
 G(A)=\log\langle\exp(q\tau(AJ))\rangle,
\]
one has
\begin{align}
 G(A)&=\frac{q^2k_{\mathrm{def}}}{2}\,
       \tau_z\tau_w\left(\frac{A(z)A(w)}{(z-w)^2}\right),
       \label{cylocal:eq:gaussian-log}\\
 sG(A)&=q^2k_{\mathrm{def}}\tau(c'A).
       \label{cylocal:eq:ward}
\end{align}
If $q^2k_{\mathrm{def}}\ne0$, this variation cannot be removed by
translation-invariant, finite-order differential-polynomial
counterterms in the source alone, with constant coefficients and
no inverse antiholomorphic derivative.
\end{proposition}
\begin{proof}
The current $J$ is linear in Gaussian fields. All connected
cumulants above order two vanish, and its second cumulant is
\eqref{cylocal:eq:defect-level}. Wick's identity gives
\eqref{cylocal:eq:gaussian-log}. The double integral is the action
of the derivative of the Cauchy distribution on smooth compact
joint sources. It involves no product of distributions on the
diagonal.

The two variations cancel the factor $1/2$. Integration by parts
and \eqref{cylocal:eq:cauchy} give
\begin{align*}
 sG(A)
 &=q^2k_{\mathrm{def}}\tau_z\tau_w
       \left(\frac{\bar\partial c(z)A(w)}{(z-w)^2}\right)\\
 &=-q^2k_{\mathrm{def}}\tau_z\tau_w
       \bigl(c(z)A(w)\partial_z\delta^\tau(z-w)\bigr)\\
 &=q^2k_{\mathrm{def}}\tau(c'A).
\end{align*}
Compact support removes the outer boundary term.

For the counterterm assertion, grade differential monomials by
their total antiholomorphic derivative order. The source gauge
differential replaces $A$ by $\bar\partial c$ and raises this order.
Holomorphic integration by parts preserves it. Antiholomorphic
total derivatives have positive order. The monomial $c'A$ has
order zero.

In its two-field component with one holomorphic
derivative, the only total-derivative relation is
$c'A+cA'=\partial(cA)$. Its class is nonzero, so it cannot be the
variation of a counterterm in the stated class.
\end{proof}

A smooth detector is explicit. With an independent odd scalar
$\eta$, take Schwartz functions
\[
 c(z)=\eta z e^{-|z|^2},\qquad A(z)=e^{-|z|^2}.
\]
Writing $t=|z|^2$ gives
\[
 \tau(c'A)=-\eta\int_0^\infty(1-t)e^{-2t}\,\mathrm dt
          =-\frac\eta4.
\]
Compact smooth cutoffs preserve nonvanishing for sufficiently
large support, by continuity of the distributional pairing.

The same coefficient can be retained algebraically. For the zeroth
normal label write $c_r=\partial^rc$, and let
$A_{\mathrm{gh}}=\Lambda(c_0,c_1,\ldots)$, with $T c_r=c_{r+1}$.
Let $B$ be the defect vertex algebra with translation $T_B$, and let
$\mathfrak m_{\mathrm{gh}}$ be the ideal generated by the ghosts.
Form the positive ghost-density
space
\[
 L_B=\frac{\mathfrak m_{\mathrm{gh}}\otimes B}
             {(T+T_B)(\mathfrak m_{\mathrm{gh}}\otimes B)}.
\]

It is completed in $\hbar,q$, with finite tangential jet support at
each precision. The vertex zero product gives a graded Lie bracket on this
quotient. Its differential is zero in the present abelian model.
For $\alpha=q\overline{c_0J}$, the residue calculation gives
\begin{equation}\label{cylocal:eq:density-curvature}
 \frac12[\alpha,\alpha]
   =\frac{q^2k_{\mathrm{def}}}{2}
         \overline{c_1c_0\mathbf1}.
\end{equation}
Indeed, the linear Taylor term of $c(z)$ multiplies the double pole
in \eqref{cylocal:eq:defect-level}.

Assume $q^2k_{\mathrm{def}}\ne0$. Take the vacuum coefficient in the chosen
normal-ordered state basis, polarize the two ghosts on Laurent
functions $f,g$, and take the residue. Total translation is killed
by this functional: $T_B$ has zero vacuum coefficient, and the
residue of a derivative vanishes. On the right side of
\eqref{cylocal:eq:density-curvature} it gives
\[
 q^2k_{\mathrm{def}}\operatorname{Res}_0(f'g\,\mathrm dz).
\]

The pair $(f,g)=(z,z^{-1})$ gives $q^2k_{\mathrm{def}}$.
The vacuum-coefficient functional used here is linear. It is not
a scalar augmentation of the vertex algebra. Since $d=0$, no
higher-order correction to $\alpha$ can cancel its order-$q^2$
curvature in $L_B$.

Descent agrees with the physical calculation. The odd derivation
$\delta c_r=\partial^r A$, $\delta A=0$, gives
\[
 \delta(c_1c_0)=A'c-c'A,
 \qquad -\tau\delta(c_1c_0)=2\tau(c'A).
\]
Applying $-\tau\delta$ to
\eqref{cylocal:eq:density-curvature} gives exactly
\eqref{cylocal:eq:ward}. For the current $\partial J$, the fourth
pole is $-6k_{\mathrm{def}}/(z-w)^4$. The Taylor factor $1/3!$
then gives curvature
$-q^2k_{\mathrm{def}}\overline{c_3c_0}/2$ and Ward functional
$-q^2k_{\mathrm{def}}\tau(c'''A)$.

There are nonzero couplings with regular OPE. Take two independent
defect systems and currents
\[
 J_+=\beta_++\frac{\kappa_{\mathrm{def}}}{2}\partial\gamma_+,
 \qquad
 J_-=\beta_- -\frac{\kappa_{\mathrm{def}}}{2}\partial\gamma_-.
\]
Their mixed contractions vanish and their levels are opposite.
Consequently $J_{\mathrm{tot}}=J_++J_-$ has regular self-OPE.
It is nonzero because the two $\beta$ coefficients are independent.

Its connected source logarithm vanishes and its Gaussian Ward
identity holds to every order. The original nonzero-level current
retains the obstruction of Proposition~\ref{cylocal:prop:ward}.

\section{The comparison targets}
\label{cylocal:sec:targets}

Consider a family of mutually regular even currents indexed by normal
labels $\alpha$. Define the free commutative differential algebra
\[
 U_0=\operatorname{Sym}
       \left(\bigoplus_{\alpha\in\mathbb N^2}\mathbb C[T]j_\alpha\right),
 \qquad Y(u,z)v=(e^{zT}u)v.
\]
It has a unital vertex map to their coefficient algebra.
This map sends $j_\alpha$ to the chosen current.

Derivatives and the vertex Wick
identity show that all generated currents remain mutually regular.
The normal product is therefore commutative and associative, which
proves this universal property. Parameter completion is taken only
after the algebraic normal-label direct sum is specified.
After scalar extension and parameter completion, the current
$J_{\mathrm{tot}}$ gives such a map by sending
$j_{(0,0)}$ to $qJ_{\mathrm{tot}}$ and the other generators to zero.

Extend the nonequivariant Hall coefficients to $\mathbb C$.
The commutator $[e_1,e_0]$ in this algebra is nonzero.
The normal product on $U_0$ is commutative. Hence no isomorphism
can carry this Hall multiplication to the normal product on $U_0$. This is an obstruction for this specified operation
comparison with the free abelian current target.

Several constructions remain available with their own operations:
\begin{center}
\small
\begin{tabularx}{\textwidth}{@{}lXX@{}}
\toprule
Object&Coefficient and operation&Unit and parameter data\\
\midrule
$\mathcal H^T\cong Y_R^+$
 &Critical Hall multiplication on all charges
 &$R=\mathbb Q[a,b]$ and charge-zero unit.\\
$\mathcal A$
 &Polynomial ambient shuffle product
 &The same ring and grading, with a map that can lose faithfulness
 after specialization.\\
$B$
 &The Gaussian defect vertex algebra and all singular products
 &Local vacuum, $\hbar,q,\kappa_{\mathrm{def}}$ and the fixed covariance.\\
$U_0$
 &Universal mutually regular current fields
 &Local vacuum, algebraic normal labels and parameter completion.\\
\bottomrule
\end{tabularx}
\end{center}

A representation on Hilbert-scheme cohomology is further data.
Already $H_T^\bullet(\operatorname{Hilb}^1(\mathbb C^2))=R$,
whereas the charge-one Hall coefficient is $R[u][2]$.
A full double additionally requires a negative half, Cartan data,
pairing and compatible completion. The positive-algebra isomorphism
of Theorem~\ref{cylocal:thm:yangian} supplies its stated associative
comparison.

\begin{problem}\label{cylocal:prob:field-comparison}
Construct a family of critical Hall coefficient complexes over
ordered configurations on a curve, with its collision and unit maps.
Compare this family with an independently specified field theory.
Specify the coefficient-ring map, the grading dictionary, and the
treatment of any coefficient dual. Construct the maps preserving
translation, vacuum and every supported collision operation.
\end{problem}

Chern roots in \eqref{cylocal:eq:shuffle} are internal equivariant
parameters. A formal generating series in those parameters is a
candidate input to this problem. A spacetime field requires its
state-field map and locality. The surviving class of
Proposition~\ref{cylocal:prop:specialization-kernel} and the explicit
Ward class of Proposition~\ref{cylocal:prop:ward} are constraints on
any proposed comparison using these carriers.
